% !TEX program = xelatex
\documentclass[12pt]{report}
\usepackage[T1]{fontenc}
\usepackage{fontspec}

\defaultfontfeatures{Ligatures=TeX}
\setmainfont{Roboto Light}
\setsansfont{Roboto Light}
\setmonofont{Roboto Mono Light}

\usepackage[english, bulgarian]{babel}
\usepackage{indentfirst}
\usepackage[a4paper, portrait, margin = 2.5 cm]{geometry}
\usepackage{url}
\usepackage{color}
\usepackage{xcolor}
\usepackage{graphicx}
\usepackage{adjustbox}
\usepackage{listings}
\renewcommand{\baselinestretch}{1.1}
\setlength{\emergencystretch}{3em}

\usepackage{graphicx}
\graphicspath{ {./pictures/} }

\lstset{
	backgroundcolor = \color{light-gray},
    language = C,
    xleftmargin = 1cm,
    framexleftmargin = 1em,
    basicstyle=\ttfamily,
	moredelim=[is][\underbar]{_}{_},
}

\usepackage{color}
\definecolor{Bluish}{rgb}{0.39,0.55,0.78}
\definecolor{light-gray}{gray}{0.9}

\usepackage{hyperref}
\hypersetup{
    colorlinks=true,
    linktoc=all,
    citecolor=black,
    filecolor=black,
    linkcolor=black,
    urlcolor=black
}

\usepackage{tabularx}


\title{Дипломна работа}
\author{Диана Генева <dageneva@qtrp.org>}
\date{2018}

\begin{document}
\maketitle
\thispagestyle{empty}
\tableofcontents
\pagebreak

\chapter{Нулева зона}
Бла, бла, бла, аз съм толкова емоционална.

\includegraphics[width=\textwidth]{pikachu}

\chapter{Емоции в реч}
    \section{Физика на тъгата}
    \section{Загладено опростяване}    
    \section{Характеристики}
        \subsection{Избор}
        \subsection{Извличане}
    \section{Класификация}    
    \section{Резултати}

\chapter{Грубо в мозъка}
    \section{Характеристики}
        \subsection{Избор}
        \subsection{Извличане}
    \section{Класификация}    
    \section{Резултати}
\chapter{Двойната звезда}
    \section{Резултати}
\chapter{Големият портрет}


\end{document}